\textbf{4. (Normandy landings).}

In 1944, the Allies invaded Normandy, which influenced the outcome of World War II.
The Allies had 6 possible attacking configurations (1--6) and the Nazis had 6 possible defensive strategies (A--F).
All possible battles were simulated and the following outcomes for the Allies were estimated
(the outcomes for the Nazis are the negations of these values):

\[
\begin{array}{c|cccccc}
 & A & B & C & D & E & F\\
\hline
1 & 13 & 29 & 8 & 12 & 16 & 23\\
2 & 18 & 22 & 21 & 22 & 29 & 31\\
3 & 18 & 22 & 31 & 31 & 27 & 37\\
4 & 11 & 22 & 12 & 21 & 21 & 26\\
5 & 18 & 16 & 19 & 14 & 19 & 28\\
6 & 23 & 22 & 19 & 23 & 30 & 34
\end{array}
\]

\begin{enumerate}
\item[(a)] Eliminate all dominated strategies.
\item[(b)] Try to solve the remaining game.
\item[(c)] In reality, the Germans chose defense \(B\) and the Allies chose attack \(1\).
Comment on these choices.
\end{enumerate}

\textbf{Solution.}


\textbf{(a) Elimination of dominated strategies.}

\textbf{Rows.}
The Allies aim to maximize.

\begin{itemize}
\item Row 4 is dominated by row 2:
\[
(11,22,12,21,21,26)\le (18,22,21,22,29,31)
\]
\item Row 5 is dominated by row 6:
\[
(18,16,19,14,19,28)\le (23,22,19,23,30,34)
\]
\end{itemize}

Thus rows \(4\) and \(5\) are eliminated.

\textbf{Columns.}
The Germans aim to minimize.

\begin{itemize}
\item Column \(E\) is dominated by \(A\):
\[
(16,29,27,30)\ge (13,18,18,23)
\]
\item Column \(F\) is dominated by \(A\):
\[
(23,31,37,34)\ge (13,18,18,23)
\]
\end{itemize}

Remove columns \(E,F\).

We obtain:
\[
\begin{array}{c|cccc}
 & A & B & C & D\\
\hline
1 & 13 & 29 & 8 & 12\\
2 & 18 & 22 & 21 & 22\\
3 & 18 & 22 & 31 & 31\\
6 & 23 & 22 & 19 & 23
\end{array}
\]

\begin{itemize}
\item Column \(D\) is dominated by \(C\):
\[
(12,22,31,23)\ge (8,21,31,19)
\]
\end{itemize}

Remove \(D\), giving:
\[
\begin{array}{c|ccc}
 & A & B & C\\
\hline
1 & 13 & 29 & 8\\
2 & 18 & 22 & 21\\
3 & 18 & 22 & 31\\
6 & 23 & 22 & 19
\end{array}
\]

\begin{itemize}
\item Row 2 is dominated by row 3:
\[
(18,22,21)\le (18,22,31)
\]
\end{itemize}

Remove row \(2\).

\textbf{Remaining game.}
\[
\boxed{
\begin{array}{c|ccc}
 & A & B & C\\
\hline
1 & 13 & 29 & 8\\
3 & 18 & 22 & 31\\
6 & 23 & 22 & 19
\end{array}}
\]

\textbf{(b) Solving the remaining game.}

\textbf{Check for pure equilibrium.}

Row minimum entries in each row:
\[
8,\quad 18,\quad 19 \quad \Rightarrow \quad \text{maximum}=19
\]

Column maximum entries in each column:
\[
23,\quad 29,\quad 31 \quad \Rightarrow \quad \text{minimum}=23
\]

There is no pure Nash equilibrium because the Nazis want to minimize their payoff, while the Allies want to maximize their payoff\\

\textbf{Mixed equilibrium.}

Let the Allies mix between rows \(3\) and \(6\) because the lowest payoff \(\ge\) than the highest payoff in row 1:

\[
p = \Pr(3), \quad 1-p = \Pr(6)
\]

Expected payoffs:

\[
u(A) = 18p + 23(1-p) = 23 - 5p
\]
\[
u(B) = 22
\]
\[
u(C) = 31p + 19(1-p) = 19 + 12p
\]

Equalizing \(A\) and \(C\):

\[
23 - 5p = 19 + 12p \;\Rightarrow\; p = \frac{4}{17}
\]

Thus the Allies play:
\[
\left(0,0,\frac{4}{17},0,0,\frac{13}{17}\right)
\]

The value is:
\[
v = 23 - 5p = 23 - 5 \cdot \frac{4}{17} = \frac{371}{17} \approx 21.82
\]

\textbf{Germans' strategy.}

Let the Germans mix between \(A\) and \(C\) because column
B gives higher payoffs to the Allies than the optimal mix and it is worse for the Germans:

\[
q = \Pr(A), \quad 1-q = \Pr(C)
\]

Indifference between rows \(3\) and \(6\):

\[
31 - 13q = 19 + 4q \;\Rightarrow\; q = \frac{12}{17}
\]

Thus the Germans play:
\[
\left(\frac{12}{17},0,\frac{5}{17},0,0,0\right)
\]

\textbf{(c) Comment on choices:}
Germany played suboptimally by choosing
B, and the Allies correctly exploited this by choosing
1, resulting in a higher payoff than the equilibrium value.
